\section{模型评价与推广}

\subsection{模型评价}

本文在同一套微网物理约束上逐步增加信息和交易规则，使四问的费用变化能够追溯到信息边界、结算方式和价格条件。统一物理内核也使各问可以共享能量与状态校验，减少因重复建模带来的定义漂移。

统一结构带来三方面优势。其一，调度结果能够直接还原为低价充电、高价放电、光伏先存后用和滚动修正等可解释动作；其二，计划信息与执行信息严格分离，并由反事实扰动检验非预见性，使费用结果对应现实时间顺序下可实施的策略；其三，线性/混合整数规划允许逐时回代能量、状态和费用，便于发现单位错误、状态越界和充放电重叠退化。

需要看到的是，模型的“严格可执行”仍建立在若干简化之上。问题二、三采用日循环终端策略，以抑制午夜前透支电池的有限时域效应，但日循环约束也限制了跨日价格和光伏差异的主动利用；历史预测器只使用已有序列，没有气象、温度、节假日等外生特征，对极端天气和异常负荷的响应有限；光伏风险边际目前按预测器使用统一残差分位，尚未刻画不同季节和日内时刻的异方差；储能也未计入自放电、容量衰减和循环寿命成本。因而，本文给出的全年费用应理解为既定储能参数和题面交易规则下的运行费用，尚未包含完整资产寿命成本。

\subsection{模型改进}
\label{sec:improve}

模型改进应优先针对上述简化。第一，可把计划时域由单日扩展到 $3\sim7$ 天，并只执行首日方案，在窗口末端加入储能残值，使模型能够利用跨日差异，并避免最后一天放空电池。第二，可引入数值天气预报、辐照度、温度、工作日等特征，建立负载与光伏的联合概率场景，再通过情景抽样或 CVaR 直接优化跨时段风险，使问题二中的单时段分位数启发转化为完整的随机调度。第三，可按季节和日内时刻分别估计光伏残差分布，使安全边际随实际预测不确定性变化，避免全天采用同一绝对修正量。第四，可加入电池循环衰减、自放电、功率相关效率和维护费用，在“少买电”和“延长储能寿命”之间建立长期经济权衡。第五，当前费用函数未计入获取日内新预报、修改交易计划所需的固定操作成本，也未考虑需求侧响应；若面向真实系统，应把这些成本与可调负荷一并纳入。

\subsection{模型推广}

本文“公共物理约束—信息分层—滚动修正”的结构具有一定的可迁移性。在园区综合能源系统中，可把单一电力平衡扩展为电、热、气多能流平衡，并增加储热、储气状态 \cite{liu2018dayahead}；在电动汽车充电站中，可把居民负载替换为车辆到站后的充电需求，再加入离站时间和最低荷电状态约束；在工业用户侧储能中，则可以在现有目标函数中加入最大需量电费和生产负荷约束。更一般地，只要系统具有“先做计划、后观察偏差、状态逐期传递、未来计划允许被新信息修正”的特征，例如库存补货、水库调度和供应链运输，都可以沿用本文的信息边界与滚动决策框架，再根据具体业务替换物理状态和费用规则。
